\section{Introduction}
\label{sec:intro}

Let \(G,\Theta\in\mathbb R^{m\times n}\), with \(\Theta\ne0\).  Our primary
object is the nuclear-norm distance from \(G\) to the one-dimensional matrix
subspace generated by \(\Theta\):
\begin{equation}
 d^\star
 =\min_{\lambda\in\mathbb R}\phi(\lambda),
 \qquad
 \phi(\lambda)=\|G+\lambda\Theta\|_*
 =\operatorname{dist}_{\|\cdot\|_*}
   \bigl(G,\operatorname{span}\{\Theta\}\bigr).
 \label{eq:D}
\end{equation}
Classical distance duality identifies the same value with
\begin{equation}
 p^\star=
 \max_{\Phi}
 \left\{\langle G,\Phi\rangle:
 \|\Phi\|_2\le1,\ \langle\Theta,\Phi\rangle=0\right\}.
 \label{eq:P}
\end{equation}
Thus \(d^\star=p^\star\): the dual unit ball is the spectral-norm ball, and
the annihilator of \(\operatorname{span}\{\Theta\}\) is one hyperplane.  We
study what happens when the affine matrix line
\(\lambda\mapsto G+\lambda\Theta\) meets the rank-deficient locus.  The
distinctive difficulty is that the natural point derivative is expressed by
the compact polar factor.  It is smooth on fixed-rank strata, but under
general rank-changing perturbations the zero-kernel completion need not have a
continuous extension; see, e.g., the classical polar-decomposition framework
of \citet{higham1986polar}.

When \(\Theta=uv^\top\), \eqref{eq:P} is the tangent LMO on the smooth
spectral sphere.  Norm-ball LMOs underlie matrix-aware first-order methods
\citep{pethick2025lmo}.  \citet{xie2026controlled} introduced the corresponding
compact-polar bisection, and \citet[Appendix P.5]{li2026intrinsic} interpreted
it as the spectral-sphere instance of an intrinsic LMO.  Neither work treats
rank-deficient crossings.

\paragraph{Classical boundary and the new object.}
Optimality in \eqref{eq:D} is classical Birkhoff--James orthogonality.
Distance duality, matrix-norm subdifferentials, and the associated criteria
are established in
\citep{singer1970best,watson1992subdifferential,zietak1993faces,
watson1993kyfan,bhatia1999orthogonality,grover2014distance,
johnston2022trace,grover2026kyfan}; we claim no novelty for those facts.
The new object is the tangency slice
\(\partial\|A^\star\|_*\cap\Theta^\perp\), its spectral centers, and the
singular asymptotics of their regularization paths.

Canonical approximant selection by Schatten limits is also classical:
\citet{legg1985canonical} treat \(p\downarrow1\), while
\citet{zietak2017strict} surveys the \(p\uparrow\infty\) strict-spectral
counterpart, including its conjectural general form.  Our selection instead
occurs on the equality-constrained certificate face.
If \(A=U_r\Sigma_rV_r^\top\) is a compact singular value decomposition,
the compact polar factor is
\begin{equation}
 \Sc(A)=U_rV_r^\top.
 \label{eq:compact-polar}
\end{equation}
At full rectangular rank, \(\phi\) is differentiable and
\(\phi'(\lambda)=\langle\Theta,\Sc(G+\lambda\Theta)\rangle\).  This suggests
the ordinary scalar equation
\begin{equation}
 h_{\rm c}(\lambda):=
 \langle\Theta,\Sc(G+\lambda\Theta)\rangle=0.
 \label{eq:compact-root}
\end{equation}
At full rectangular rank this equation is valid.  When the affine matrix loses
rank, it need not remain valid, because the complete subdifferential is
\begin{equation}
 \partial\|A\|_*
 =\left\{
 U_rV_r^\top+U_0KV_0^\top:\ \|K\|_2\le1
 \right\},
 \label{eq:full-subdiff}
\end{equation}
where \(U_0,V_0\) span the left and right null spaces.  The compact factor is
only the choice \(K=0\).  The omitted kernel block can be exactly what is
needed to enforce \(\langle\Theta,\Phi\rangle=0\).  Consequently,
\(h_{\rm c}\) can be monotone and still jump across zero.

Thus the LMO and nuclear-norm duality remain exact; only the replacement of a
set-valued optimality graph by one compact selection can fail.

\subsection{Contributions and novelty boundary}

The contributions, in the order used by the analysis, are as follows.

\begin{enumerate}
 \item Starting from the classical interval
 \(\partial\phi=[a-\beta,a+\beta]\)
 \citep{watson1992subdifferential,lewis1995convex}, we classify feasibility
 and uniqueness of the tangency slice.  Its minimum-Frobenius member follows
 by singular-value water-filling; clipping is universally unnecessary exactly
 for scalar multiples of partial isometries.

 \item An admissible Fenchel family unifies the hard and smooth centers:
 Huber gives \(K_F\), pseudo-Huber gives \(K_R\), and we characterize their
 agreement.  A polynomial derivative tail of order \(p>1\) yields exponent
 \(p/(p+1)\), while finite saturation gives a linear law.  The novelty beyond
 standard smoothing principles
 \citep{nesterov2005smooth,lobos2025smooth,attouch1984variational} is the
 constrained spectral center and its coupling to the singular multiplier.

 \item We prove the regular expansion and, at every strict rank-deficient
 minimizer, analytic multiplier and returned-direction expansions without
 shape, normal-rank, or corank restrictions.  Compact-root instances have
 computable cubic/quadratic superconvergence and a quintic/quartic refinement.
 At transverse square corank-one boundaries, pseudo-Huber has the exact
 two-thirds law and a uniform cubic strict-to-boundary crossover.  Our new
 content is the matrix-line exponents, constants, and crossover, rather than
 the existence of fractional central-path rates \citep{basu2022central}.

 \item Finally, a strict \(2\times2\) witness extends to a nonempty open
 compact-polar failure set in every square dimension at least two, even with
 unique approximant and certificate.  This motivates retaining the full
 slice.  Rank-one normals also admit an SVD-free crossing formula.
\end{enumerate}

An open failure set implies neither an event-frequency law nor an outer
algorithmic or learning-performance gain; those require additional models.

\subsection{Organization}

\Cref{sec:duality,sec:failure} give the scalar interval, structural failure,
and certificate geometry; \cref{sec:smoothing,sec:rates} develop the Fenchel
centers and local laws; \cref{sec:numerics} verifies the constants.  Auxiliary
level-set geometry and certificate repair are collected in the supplement.
